\documentclass[11pt]{article}

\usepackage[margin=1in]{geometry}
\usepackage{amsmath,amssymb,amsthm,mathtools}
\usepackage{booktabs,tabularx,array}
\usepackage{pdflscape}
\usepackage{float}
\usepackage{microtype}
\usepackage{enumitem}
\usepackage{xcolor}
\usepackage{xurl}
\usepackage[hidelinks]{hyperref}

\definecolor{verified}{RGB}{0,92,74}
\definecolor{external}{RGB}{45,76,135}
\definecolor{conjectural}{RGB}{145,76,0}
\newcommand{\Mtag}{\textcolor{verified}{\textsc{M}}}
\newcommand{\Ttag}{\textcolor{external}{\textsc{T}}}
\newcommand{\Ctag}{\textcolor{conjectural}{\textsc{C}}}
\newcommand{\dd}{\mathrm{d}}
\newcommand{\ii}{\mathrm{i}}
\newcommand{\tr}{\operatorname{tr}}
\newcommand{\rank}{\operatorname{rank}}
\newcommand{\Span}{\operatorname{span}}
\newcommand{\diag}{\operatorname{diag}}
\newcommand{\R}{\mathbb{R}}
\newcommand{\Q}{\mathbb{Q}}
\newcommand{\N}{\mathbb{N}}
\newcommand{\Z}{\mathbb{Z}}
\newcommand{\C}{\mathbb{C}}
\newcommand{\Mass}{\mathcal{M}}

\newtheorem{theorem}{Theorem}[section]
\newtheorem{proposition}[theorem]{Proposition}
\newtheorem{corollary}[theorem]{Corollary}
\newtheorem{definition}[theorem]{Definition}
\newtheorem{remark}[theorem]{Remark}
\newtheorem{conjecture}[theorem]{Conjecture}

\title{Mass as Null-Edge Spectral Area:\\
Machine-Verified Null-Edge Kinematics and a Finite Carrier Model}
\author{Null-Edge Formalization Project}
\date{July 9, 2026}

\begin{document}
\maketitle

\begin{abstract}
We isolate a finite mathematical core for the statement that massive momentum
can be assembled from massless constituents.  If spinors
$\psi_1,\ldots,\psi_N\in\C^2$ define the positive semidefinite momentum Gram
matrix $P=\sum_i\psi_i\psi_i^\dagger$, then
\[
  \det P=\sum_{i<j}|\psi_i\wedge\psi_j|^2.
\]
Rank distinguishes the massless boundary from the massive interior: the rank
is one exactly when all occupied null directions are collinear.  The continuous
mass squared is not the rank itself but the determinant, the summed spinor area.
This is classical spinor kinematics, formalized in Lean 4.  We then study a
finite Krein-space Dirac carrier whose adjoint square
separates into aperture, closure, turn, and soldering-shaped blocks.  The
explicit rational four-block witness is coefficient-rigid, although the same
split is not forced by grading type alone.  A positive physical sector, an
  exact aperture--closure phase diagram, a finite interacting bound state, and
  an exact $1+1$ checkerboard path sum are also machine verified.  In the latter,
  each null history carries weight $(\ii\epsilon m)^{\#\text{corners}}$ and the
  finite sum satisfies a discrete Dirac recursion.  Its one-step symbol obeys
  the exact lattice relation $\cos(\omega a)=\cos(ka)\cos(ma)$; at zero momentum
  its eigenphases are $e^{\mp\ii ma}$, so the turn parameter is the
  principal-branch spectral gap.  Further kernel checks give a
  $3+1$ tetrahedral endpoint law, ordered Weyl-projector path amplitudes, and a
  normalization obstruction,
  a concrete quotient $D_4/L_H\simeq(\Z/2)^3$ of cardinality eight,
  a common null-characteristic theorem for lower-order channel mixing, a
nontrivial gauge-invariant finite closure holonomy, and a nondegenerate
rational matrix coframe.  We distinguish these finite
theorems from the conjectural identifications of the four blocks with kinetic,
QCD, Yukawa, and gravitational mass in a continuum theory.  The result is a
falsifiable finite model and proof architecture, not a derivation of measured
particle masses or the Standard Model.
\end{abstract}

\paragraph{Claim labels.}
Every substantive statement is marked as external mathematics (\Ttag),
machine-verified in the pinned Lean project (\Mtag), or conjectural (\Ctag).
The labels are part of the scientific claim: a formal proof verifies the Lean
statement, while the physical interpretation remains a separate question.

\section{Introduction}

A future-directed timelike four-momentum can be written as a sum of
future-directed null momenta.  In spinor notation this familiar fact takes a
particularly economical form.  A null momentum is a rank-one positive
semidefinite bispinor, and a sum of null momenta is a positive semidefinite
$2\times2$ matrix.  The determinant measures the failure of that sum to remain
rank one.  In standard normalization it is the invariant mass squared.

The first purpose of this paper is to make this kinematic statement exact,
finite, and machine checked.  The second is to ask how much of a dynamical
operator model can be built while retaining the same discipline.  We introduce
a finite Dirac-type carrier
\begin{equation}\label{eq:carrier}
  D=\sum_e c(\alpha_e)\nabla_e+\Gamma\phi,
\end{equation}
where $c(\alpha_e)^2=0$ encodes null edge directions, $\nabla_e$ is finite
transport, and $\Gamma\phi$ changes the chiral or turn sector.  Its
Krein-adjoint square has four algebraically distinguished contributions.  The
neutral names used below are aperture, closure, turn, and soldering gradient.

The central hierarchy is:
\begin{enumerate}[label=(\roman*)]
  \item the Gram determinant is the theorem-bearing mass invariant;
  \item the finite carrier is a candidate dynamical realization of that
        invariant;
  \item the correspondence between carrier blocks and continuum Standard Model
        terms is a testable conjecture.
\end{enumerate}
Conflating these levels would make the model appear stronger while making it
scientifically weaker.  We therefore state positive results beside their
sharpest known obstructions.

The main finite contributions are:
\begin{itemize}[leftmargin=1.5em]
  \item an exact Cauchy--Binet mass identity and its massless iff collinear
        boundary (\Ttag/\Mtag);
  \item a uniqueness theorem for the determinant among quadratic forms
        vanishing on all null edges (\Mtag);
  \item exact coherence, entropy, and path-decoherence forms of the same
        two-direction invariant (\Mtag);
  \item the Mandelstam/spinor-helicity dictionary, termwise
        $\mathrm{SL}(2,\C)$ invariance, two-edge concurrence identity, and
        square-root mass superadditivity on the future cone (\Ttag/\Mtag);
  \item a finite carrier-square decomposition, together with a concrete
        coefficient-rigidity theorem and an abstract non-uniqueness result
        (\Mtag);
  \item a positive-sector mass witness and a complete finite
        aperture--closure phase diagram (\Mtag);
  \item an exact Gaussian-rational checkerboard path sum satisfying a discrete
        Dirac recursion, together with its exact quasienergy relation and
        zero-momentum eigenphases (\Mtag);
  \item exact $3+1$ tetrahedral endpoint kinematics, ordered Weyl-projector
        bend/orientation phases, a sharp normalization obstruction, and a
        conditional common null-front theorem (\Mtag);
  \item a nontrivial history-local $U(1)$ closure holonomy and a nondegenerate
        finite matrix coframe with invariant Lorentzian defect action (\Mtag);
  \item a finite second-quantized gap and a derived two-body closure interaction
        with both binding and no-binding regimes (\Mtag).
\end{itemize}

Prior work already relates massive momenta to pairs of null spinors, momentum
bispinors to two-qubit entanglement, and Dirac propagation to checkerboard or
quantum-walk histories \cite{ElvangHuang,ArkaniHamedMassive,ChinLee,FeynmanHibbs,JacobsonSchulman}.
Our contribution is not the physical picture itself.  It is the finite
formalization, the integration with a single carrier square, and a proof policy
that records nondegenerate witnesses and kill conditions beside headline
claims.

\subsection*{Position relative to discrete relativistic dynamics}

Quantum walks and quantum cellular automata already derive free Weyl and Dirac
evolution from homogeneity, locality, unitarity, isotropy, or information-
processing principles, in both one and three spatial dimensions
\cite{DarianoPerinottiPrinciples,BisioDarianoTosini,DarianoPathIntegral,
DarianoWeyl,MlodinowBrun}.  Their continuum limits and symmetry classifications
are substantially more developed than ours.  The present contribution is
different: it connects an exact corner-weighted history sum to the same
determinant/rank invariant used by a four-channel finite carrier, and exposes
the resulting statements through kernel-checked declarations and explicit
failure modes.  It is a formal bridge into that literature, not the first
discrete derivation of the Dirac equation.

Discrete gravitational work likewise contains far richer geometry: Regge
calculus places curvature on simplicial hinges; connection formulations encode
curvature by holonomy; twisted and simplicial phase spaces impose closure and
gluing constraints; teleparallel gravity uses a nondegenerate coframe and flat
connection; and position-dependent quantum walks model curved backgrounds
\cite{Regge,DittrichRyan,FreidelSpeziale,BaezWise,ArrighiCurvedWalk}.  Our
matrix-coframe result proves only the finite local-frame, metric, volume,
refinement, and nonzero-defect gates needed before a continuum comparison is
meaningful.  It does not compete with those continuum or simplicial dynamics.

\section{The candidate theory}
\label{sec:theory}

This paper defends one identifiable proposal.  We state it here in full, with
each load-bearing arrow graded, so that the finite results of the later
sections read as evidence for a single architecture rather than as a
collection of related lemmas.

\subsection{Postulates}

\begin{enumerate}[label=\textbf{P\arabic*.},leftmargin=2.2em]
  \item \textbf{Null primitives.}  The primitive data are finite oriented null
  events: rank-one lightlike transmissions carrying spinor direction, internal
  (charge, chirality, strand) labels, and complex amplitude.  There is no
  primitive timelike motion, no primitive mass, and no background spacetime
  container.
  \item \textbf{Coherent composition.}  Histories compose by gluing and
  superpose linearly; amplitudes of composite histories are ordered products
  of elementary transport, turn, and soldering gates.  (The finite gluing
  algebra is \Mtag; the claim that nature's amplitudes are exhausted by such
  finite sums is \Ctag.)
  \item \textbf{Gauge redundancy.}  Descriptions related by the constraint
  differential $Q$ ($Q^2=0$) are identical physics.  Physical content lives on
  the cohomology $\ker Q/\mathrm{im}\,Q$ (finite Kugo--Ojima structure,
  \Mtag).
  \item \textbf{Positive decoding.}  The physical sector is the
  Krein-positive part of cohomology; probabilistic quantum theory operates
  there and only there.  (Positive-sector existence and its sign no-go are
  \Mtag; the normalized outcome rule itself is at present an imported Born
  input, not a derivation.)
  \item \textbf{One carrier square.}  Time evolution and interaction enter
  through one finite Dirac-type carrier $D$; every second-order physical
  effect is one of the four graded blocks of its Krein square --- aperture,
  closure, turn, soldering gradient.  (The four-type decomposition is forced,
  \Mtag; uniqueness of the split is not, and the equivalence class of
  presentations modulo chain homotopy and Krein intertwiners is the physical
  object, \Mtag.)
\end{enumerate}

\subsection{Derivation spine}

\begin{center}
\small
\begin{tabularx}{\linewidth}{@{}>{\raggedright\arraybackslash}p{0.30\linewidth}>{\raggedright\arraybackslash}p{0.12\linewidth}X@{}}
\toprule
Arrow & Grade & Anchor / open bridge\\
\midrule
null histories $\to$ gauge classes & \Mtag & finite Kugo--Ojima; generic
Hilbert--Hodge representatives\\
gauge classes $\to$ positive states & \Mtag/\Ctag & positive Hodge decoder;
invariant sector selection open\\
positive states $\to$ mass & \Mtag & Gram determinant $=$ Cauchy--Binet wedge
sum; one supplied pair gives the same Hodge cost and finite-action Hessian\\
mass $\to$ dynamics & \Mtag/\Ctag & exact checkerboard sum, quasienergy
relation, one-step $O(\epsilon^2)$ bound; action-derived Euler--Lagrange flow
with a positive-definite conserved first integral; unitary successive-axis
split-step walk with relativistic square and tangent generator $-iH$;
explicit sign-basis conjugacy gives the finite spatial tangent; compact-box
$3+1$ rate and the physical position/PDE limit open; a generic positive
least-residual path action has zero locus exactly equal to any selected
unitary link equation\\
states $\to$ outcomes & \Mtag/import & finite instrument API: normalized
nonnegative outcome probabilities, projective repeatability, compatible
no-disturbance; the Born trace rule itself is an imported input (P4)\\
dynamics $\to$ interactions & \Mtag/\Ctag & four-block square, coefficient
rigidity, signed closure binding; continuum field dictionary open\\
interactions $\to$ geometry & \Mtag/\Ctag & nondegenerate matrix coframe,
chain spectral distance $=$ weighted geodesic; continuum tetrad open\\
geometry $\to$ cosmology & \Mtag/\Ctag & event-count/$\Lambda$ arithmetic,
geometry-register coherence, ensemble vacuum-shift redundancy; volume
statistics imported\\
theory $\to$ experiment & V0--V2 & exact benchmark suite
(\texttt{Scripts/sim/null\_information\_lab.py}: twenty families),
including three disclosed V2 reproductions (fixed-momentum free $1+1$ and
$3+1$ Dirac propagation; independent three-level fluctuation response);
no V4 prediction yet\\
\bottomrule
\end{tabularx}
\end{center}

The finite center of this spine is now two composable machine-checked theorems
sharing one supplied pair and one literal mass parameter, not a verbal
identification.  For one supplied noncollinear spinor pair, the same
Pluecker scalar is its Hodge decoder cost, its action curvature, the stiffness
of its conserved variational flow, its Gibbs gap, and the mass parameter in
the internal $3+1$ Dirac dispersion and split-step tangent.  The rational
fixture has $m=4/25$ and
$H(1,2,2,m)^2=(5641/625)I$; a collinear pair collapses the mass, action
curvature, fluctuation, and zero-momentum Dirac generator.  This establishes a
coherent finite candidate architecture.  It does not yet prove that the scalar
recurrence and spinor walk descend from one field-valued action, nor that the
finite symbol has the required physical continuum limit.  The two theorems
therefore close parameter identity across the clusters, not identity of their
state carriers or actions.

\subsection{What the theory says}

In this theory a particle is a stable Krein-positive cohomology class of the
carrier; its mass is the least spectral cost of decoding that class, which at
the kinematic level is exactly the determinant of its null-direction Gram
matrix.  Interactions are the four ways the carrier square can fail to be
free: directional aperture, loop closure, chiral turn, and soldering drift.
The finite construction encodes spacetime information algebraically --- the
Pauli map, the null cone, and the direction sets are primitive inputs ---
while causal order and a weighted geodesic distance are reconstructed from
the carrier (finitely, \Mtag); emergence of a \emph{dynamical} continuum
geometry remains conjectural (\Ctag).  The cosmological constant is the
variable conjugate to total event
count; it is observable exactly through geometry-register coherence between
branches of different count (\Mtag, finitely), and a uniform vacuum offset is
a redundancy of the ensemble description rather than a physical energy
(\Mtag\ for the extensive case, with the exact non-extensive negative
control landed).

\paragraph{The power ledger.}  One mass operator generates several related
quantities, and the paper distinguishes their operator orders throughout:

\begin{center}
\small
\begin{tabular}{@{}lr@{}}
\toprule
Layer & Natural quantity\\
\midrule
null-spinor kinematics & $\mu^2=\det P$\\
$D^{\#}D$, action Hessians, quadratic costs & $\mu^2$\\
Dirac symbol eigenvalue / rest gap & $\mu$\\
zero-momentum quasienergy & $\mu$, modulo the lattice branch\\
exact unitary corner ratio & $\tan(a\mu)$, not $a\mu$, at finite spacing\\
\bottomrule
\end{tabular}
\end{center}

Squares give quadratic costs, absolute values give spectral gaps, and
exponentials give finite-time evolution --- of one operator, not of several
separately introduced parameters.  The construction realizing this
(the Pl\"ucker-derived mass operator $B_z$ with $B_z^2=\det P\,I$, and the
exact $\tan(a\mu)$ corner--walk correspondence) is the subject of the
in-flight formalization pair recorded in the verification section; each
downstream identification is labeled \emph{forced}, \emph{derived},
\emph{realized}, or \emph{conjectured} accordingly.

\subsection{What the theory does not yet derive}

Nine bridges are open.  Each is a named theorem target with an exact
statement, not a vague possibility: (i) the \emph{many-step continuum limit}
on a compact $3+1$ momentum domain followed by a position/PDE lift;
(ii) a \emph{shared field-valued action} whose reductions yield both the
scalar variational recurrence and the spinor Dirac walk; (iii) an
\emph{invariant positive-sector selection} principle replacing the current
witness-level choice; (iv) a \emph{local net}: tensor microcausality and
no-signaling beyond two regions; (v) the \emph{Born rule}, currently imported
at P4; (vi) the \emph{Higgs self-mass} and the absolute mass scale, for which the
theory presently offers mechanisms only for ratios and protected smallness
(dimensional transmutation is the named route to the scale);
(vii) the \emph{gauge group}: deriving the local automorphism group and its
representations rather than supplying charge registers; (viii) the
\emph{initial state and total event count}, on which the theory is silent;
and (ix) the \emph{generation count}, where the landed no-go shows a
rank-fixing input is required and the monodromy dichotomy confines the
remaining mechanism space.  The
falsifier list in Section~\ref{sec:falsifiers} states what failure looks like
for each supported arrow.  The benchmark suite executes the finite content of
the arrows assigned live S-rows in the evidence matrix; an \Mtag without such
a row remains theorem-tested rather than simulation-tested.

\section{The canonical null-edge mass invariant}

\subsection{Positive spinor Gram matrices}

Let $\psi_i\in\C^2$.  Write
\begin{equation}\label{eq:gram}
  M=\begin{bmatrix}\psi_1&\cdots&\psi_N\end{bmatrix},
  \qquad P=MM^\dagger=\sum_i\psi_i\psi_i^\dagger.
\end{equation}
Each summand is positive semidefinite and rank at most one.  Under the Pauli
map, it is a future null momentum.  The determinant of $P$ is the Lorentzian
quadratic invariant of the total momentum.

\begin{theorem}[Null-edge Cauchy--Binet, \Mtag]\label{thm:cb}
For the matrix in \eqref{eq:gram},
\begin{equation}\label{eq:cb}
  \det P=\sum_{i<j}
  \left|\det\begin{bmatrix}\psi_i&\psi_j\end{bmatrix}\right|^2.
\end{equation}
In particular $\det P\geq0$, and $\det P=0$ iff the columns of $M$ span a
subspace of dimension at most one.
\end{theorem}

Equation \eqref{eq:cb} is the exact sense in which mass is pairwise disagreement
of null directions.  The wedge vanishes when two directions agree; its squared
magnitude is the opened spinor area.  No cancellation between different pairs
is possible because every term is nonnegative.
Nonnegative scalar weights can be included by replacing each column by
$\sqrt{w_i}\psi_i$; this weighted form is a classical corollary, while
Theorem~\ref{thm:cb} is the direct Lean declaration cited below.

\paragraph{Mandelstam dictionary.}
Let $p_i=\psi_i\psi_i^\dagger$ be the null momentum associated with the $i$th
spinor and use the mostly-minus Pauli-map normalization.  Then
\begin{equation}\label{eq:mandelstam}
 |\psi_i\wedge\psi_j|^2
   =2p_i\mathbin{\cdot}p_j
   =(p_i+p_j)^2=:s_{ij},
 \qquad
 P^2=\sum_{i<j}s_{ij}.
\end{equation}
Thus Theorem~\ref{thm:cb} is exactly the familiar
$\langle ij\rangle[ji]$ or pairwise-Mandelstam expansion
\cite{ElvangHuang,ArkaniHamedMassive}.  We state this explicitly to separate
what is classical from what is contributed here: the new payload is the finite
formalization, the null-boundary uniqueness theorem, and the integration with
the carrier and history layers.  Moreover
\begin{equation}\label{eq:termwiseSL2}
 (g\psi_i)\wedge(g\psi_j)=\det(g)(\psi_i\wedge\psi_j),
 \qquad g\in\mathrm{SL}(2,\C),
\end{equation}
so every summand in \eqref{eq:cb}, not merely their total, is individually
proper-orthochronous Lorentz invariant (\Mtag for the algebraic
$\mathrm{SL}(2,\C)$ statement).

\begin{corollary}[Massive iff rank two, \Ttag/\Mtag]
For a nonzero positive semidefinite $2\times2$ momentum matrix,
\[
  \det P=0\iff\rank P=1,
  \qquad
  \det P>0\iff P\text{ is positive definite}.
\]
Thus the massless locus is precisely the rank-one boundary of the future cone.
\end{corollary}

The converse is equally important.  Every positive timelike momentum matrix has
a spectral decomposition into two positive rank-one matrices, hence every
future timelike four-momentum is a sum of future null momenta (\Mtag).  The
decomposition is not unique; the invariant is.

\subsection{Why the determinant is not an arbitrary choice}

One might object that the determinant was selected because it gives the desired
answer.  The finite formalization supports a stronger statement.

\begin{theorem}[Quadratic uniqueness, \Mtag]\label{thm:detunique}
Let $Q$ be a quadratic form on symmetric $2\times2$ matrices.  If $Q$ vanishes
on every rank-one null-edge matrix, then $Q$ is a scalar multiple of the
determinant.
\end{theorem}

The proof is coefficient-level polynomial algebra.  An explicit control,
$Q'(\left[\begin{smallmatrix}a&b\\b&c\end{smallmatrix}\right])=ac$, vanishes on
the two coordinate null rays but not on every null ray.  Universal
null-vanishing is therefore load-bearing.  Congruence invariance under
$\mathrm{SL}(2)$ is satisfied by the determinant but is not needed to force the
quadratic form in this dimension.

\subsection{A free operator bridge and an interacting correction}

The Gram invariant and the carrier spectrum are not two definitions of mass.
In the free finite model they meet through
\begin{equation}\label{eq:adjbridge}
  P\,\operatorname{adj}(P)=\det(P)I,
\end{equation}
so the free mass operator has least eigenvalue $\det P$ (\Mtag).  Interaction
changes that spectral value.

The free kinematic baseline is itself sharp.  If $p$ and $q$ are
future-causal momenta and $m(p)=\sqrt{p^2}$, the specialized Minkowski
determinant inequality gives
\begin{equation}\label{eq:superadditive}
 m(p+q)\geq m(p)+m(q)
 \qquad(\Mtag).
\end{equation}
The landed theorem proves the square-root inequality on the concrete future
cone; the classical equality characterization by proportional future-causal
momenta is not yet part of that declaration.  Equation
\eqref{eq:superadditive} is the correct free foil for binding: a bound-state
ground energy below the free constituent threshold requires an interaction,
not a failure of relativistic momentum algebra.

For the normalized coupled block used below,
the exact binding identity is
\begin{equation}\label{eq:bindinginfo}
  G_{\mathrm{bind}}=\kappa=C(\rho)\lambda,
  \qquad
  m_{\mathrm{ground}}=m_{\mathrm{free}}-C(\rho)\lambda,
\end{equation}
with signed defect $\Delta=-\kappa$ (\Mtag).  This is the expected behavior:
binding mass is not additive.  It lowers a spectral ground state relative to
the free threshold while leaving \eqref{eq:superadditive} intact for freely
added four-momenta.  Equation \eqref{eq:bindinginfo} is a finite carrier
identity, not a derivation of a QCD hadron mass.

\section{Coherence and retained direction information}

Normalize a nonzero Gram matrix by $\rho=P/\tr P$.  In two dimensions,
\begin{equation}\label{eq:linent}
  S_{\mathrm L}(\rho)=1-\tr(\rho^2)=2\det\rho.
\end{equation}
Thus linear entropy and normalized mass squared are the same scalar up to a
factor of two.  This does not identify all information measures with mass; it
identifies one exact two-level invariant.

\subsection{The two-edge concurrence identity}

For exactly two occupied edges, assemble the spinors as the columns of
$M\in\C^{2\times2}$ and read the same coefficient matrix as an unnormalized
two-qubit pure state.  Its pure-state Wootters concurrence is
$C_W(M)=2|\det M|$ \cite{Wootters}.  Since $P=MM^\dagger$,
\begin{equation}\label{eq:concurrence}
 4\det P=C_W(M)^2,
 \qquad
 \det P=0\Longleftrightarrow C_W(M)=0
 \qquad(\Mtag).
\end{equation}
For a normalized two-qubit state this is the standard concurrence; without
normalization it is the corresponding homogeneous algebraic identity.  It does
not say that a particle's rest mass is ordinary laboratory entanglement.  It
says that the two-edge Pl\"ucker invariant and the pure-state concurrence are
the same determinant read in two established dictionaries.

\subsection{Two-history complementarity}

For two visible null directions with maximal disagreement scale $D$, the
path-decoherence model gives
\begin{equation}\label{eq:pathmass}
  \det\rho(t)=t(2-t)D,
  \qquad 0\leq t\leq1.
\end{equation}
Writing hidden overlap $k=1-t$ turns this into
\begin{equation}\label{eq:duality}
  \frac{\Mass^2_{\mathrm{vis}}}{\Mass^2_{\max}}+|k|^2=1.
\end{equation}
Both forms and the rational witness $16/25+9/25=1$ are machine verified.  The
state spaces in the two derivations remain distinct; what is identified is the
normalized complementarity curve.

\subsection{Decoherence monotonicity and its boundary}

For the real trace-one block
\[
  \rho(p,x)=\begin{bmatrix}p&x\\x&1-p\end{bmatrix}
\]
and the pinching map $\Phi_t$ that sends $x\mapsto(1-t)x$, direct calculation
gives
\begin{equation}\label{eq:pinch}
  \det\Phi_t(\rho)=\det\rho+(2t-t^2)x^2.
\end{equation}
Hence pinching cannot decrease determinant on $0\leq t\leq1$ (\Mtag).  Full
pinching sends the massless coherent state
$\rho(1/2,1/2)$ to the maximally mixed state with determinant $1/4$.

The allowed-operation class matters.  A coherent signed closure rotation before
pinching can lower the final determinant relative to naive pinching (\Mtag).
Closure is therefore not interchangeable with noise.  The verified resource
statement is monotonicity for the displayed coarse-graining class, not for all
finite dynamics.

\section{A finite Dirac carrier and its square}

\subsection{The Krein-square identity}

Let $\#$ denote the Krein adjoint.  Under explicit Clifford, transport, and
covariantly constant turn hypotheses, the carrier \eqref{eq:carrier} satisfies
\begin{align}\label{eq:master}
  4D^\#D &= Q_A^\#+Q_C^\#+4Q_T+4E_\#,\\
  Q_A^\#&=\sum_{e,f}g(e,f)
    (\nabla_e^\#\nabla_f+\nabla_f^\#\nabla_e),\nonumber\\
  Q_C^\#&=\sum_{e,f}(\gamma_e\gamma_f-\gamma_f\gamma_e)
    (\nabla_e^\#\nabla_f-\nabla_f^\#\nabla_e),\nonumber\\
  Q_T&=\phi^2,\nonumber\\
  E_\#&=\sum_e\gamma_e\Gamma\phi(\nabla_e^\#-\nabla_e).\nonumber
\end{align}
This is a finite algebraic identity (\Mtag).  The terms are defined
independently; $E_\#$ is not a residual introduced to make the sum true.

The names encode operator shapes:
\begin{center}
\begin{tabularx}{\textwidth}{@{}l >{\raggedright\arraybackslash}p{0.19\textwidth} X >{\raggedright\arraybackslash}p{0.24\textwidth}@{}}
\toprule
Block & Neutral role & Continuum analogy & Verified sign information\\
\midrule
$Q_A$ & symmetric transport, aperture & covariant kinetic term & positive in the explicit sector\\
$Q_C$ & commutator, closure & curvature or chromomagnetic term & signed; balanced on the closure quotient\\
$Q_T$ & turn field squared & scalar/Yukawa term & sign fixed by the turn convention\\
$E$ & varying soldering defect & tetrad, torsion, nonmetricity & finite geometric split only\\
\bottomrule
\end{tabularx}
\end{center}
The continuum analogies are \Ctag.  In particular, the displayed theorem uses a
covariantly constant $\phi$, so its turn term is indistinguishable from an
explicit Dirac mass.  A varying-soldering theorem produces a separate defect
$E$; identifying $E$ with $E_\#$ is not proved.

\subsection{A finite local-frame gate for soldering}

The first geometric transformation-law test for the soldering interpretation
can be stated without a continuum manifold.  Represent the coframe data at the
ends of an edge $x\to y$ by rational vectors $e_x,e_y$ and its transport by
$U_{xy}$.  Define the finite defect
\begin{equation}\label{eq:solderdefect}
  T_{xy}=e_y-U_{xy}e_x.
\end{equation}
Under independent endpoint frame changes
$e_x'=g_xe_x$, $e_y'=g_ye_y$, and
$U_{xy}'=g_yU_{xy}g_x^{-1}$, the kernel proves
\begin{equation}\label{eq:soldercov}
  T_{xy}'=g_yT_{xy}.
\end{equation}
It also proves the exact refinement identity
\begin{equation}\label{eq:solderrefine}
  T_{xz}=T_{yz}+U_{yz}T_{xy},
  \qquad U_{xz}=U_{yz}U_{xy}.
\end{equation}
Consequently $\lVert T_{xy}\rVert^2$ is invariant under orthogonal endpoint
frames.  Closed-loop holonomy defects transform covariantly, while
$\tr H$ is invariant under conjugation (\Mtag).  An exact $2\times2$ rational
fixture has $T_{xy}=(0,1)\neq0$, a nonidentity orthogonal frame change, and
$\lVert T_{xy}\rVert^2=1$ (\Mtag).

The matrix-coframe upgrade now closes the finite nondegeneracy gate.  For a
square coframe $e$ and internal metric $\eta$, define
\begin{equation}\label{eq:inducedmetric}
  g[e]=e^T\eta e,\qquad \operatorname{vol}(e)=\det e,
  \qquad {\cal S}_E(T)=\tr(T^T\eta T).
\end{equation}
Invertible frame changes preserve $\det e\ne0$; an $\eta$-Lorentz frame
preserves $g[e]$ and ${\cal S}_E$, and a determinant-one frame preserves
$\operatorname{vol}(e)$ (\Mtag).  Matrix defects obey the same covariance and
refinement laws \eqref{eq:soldercov}--\eqref{eq:solderrefine}.  The exact
rational witness uses
\begin{equation}\label{eq:rationalboost}
 \eta=\begin{pmatrix}1&0\\0&-1\end{pmatrix},\qquad
 \Lambda=\begin{pmatrix}5/3&4/3\\4/3&5/3\end{pmatrix},
 \qquad \Lambda^T\eta\Lambda=\eta,\quad\det\Lambda=1,
\end{equation}
with two nondegenerate coframes and nonzero defect
$T=\operatorname{diag}(1,0)$ satisfying ${\cal S}_E(T)=1$ before and after the
nonidentity boost (\Mtag).

This passes finite covariance, nondegeneracy, metric, volume, action, and
subdivision gates suggested by discrete teleparallel geometry
\cite{BaezWise}.  It still supplies only a rational finite coframe: continuum
local Lorentz covariance, a nondegenerate tetrad field on a manifold, and a
gravitational field equation are not consequences of these identities.

\subsection{What is rigid, and what is not}

The grading types force a four-block expansion with no fifth type (\Mtag), but
they do not force a unique carrier.  Finite counterexamples admit inequivalent
complements with the same coarse type data (\Mtag).  This negative result is not
in tension with the following concrete theorem.

\begin{theorem}[Concrete coefficient rigidity, \Mtag]\label{thm:rigid}
For the explicit rational matrices $Q_A,Q_C,Q_T,E_s$ used by the carrier
witness, there are rational entry functionals $r_A,r_C,r_T,r_E$ such that
\[
 r_X(aQ_A+cQ_C+tQ_T+eE_s)=
 \begin{cases}a&X=A,\\c&X=C,\\t&X=T,\\e&X=E.\end{cases}
\]
Consequently the four matrices are linearly independent, every presentation
has unique coefficients, and $4D^TD$ recovers $(1,1,1,1)$.
\end{theorem}

The entry selectors use concrete support data.  They are exactly the additional
structure absent from the abstract non-uniqueness theorem.  The strongest
honest statement is therefore: the witness decomposition is coefficient-rigid;
the general carrier category is not yet rigid.

\subsection{Positivity and the aperture--closure phase diagram}

Krein self-adjointness does not imply positivity.  The project therefore
constructs an explicit $12$-dimensional Clifford--color carrier and compresses
it to a six-dimensional positive sector.  The compressed form is a pair of
$3\times3$ blocks.  The basic block is
\begin{equation}\label{eq:B}
 B(\lambda,\kappa)=
 \begin{bmatrix}
 \lambda&\ii\kappa&0\\
 -\ii\kappa&\lambda&0\\
 0&0&\lambda
 \end{bmatrix},
 \qquad
 \operatorname{spec}B=\{\lambda-\kappa,\lambda,\lambda+\kappa\}.
\end{equation}

\begin{proposition}[Exact finite phase diagram, \Mtag]
For $\lambda>0$,
\begin{enumerate}[label=(\alph*)]
  \item $B$ is positive definite iff $|\kappa|<\lambda$;
  \item $B$ is singular on $|\kappa|=\lambda$;
  \item for $\kappa>\lambda$, $\lambda-\kappa<0$ is an unstable direction.
\end{enumerate}
For $0\leq\kappa\leq\lambda$, the least eigenvalue is
$\lambda-\kappa$.
\end{proposition}

At $(\lambda,\kappa)=(2,1)$ the actual carrier compression is exactly the
block pair $B(2,1)\oplus B(2,-1)$ (\Mtag).  Extending this exact carrier
identification to the whole parameter plane remains open; the spectral theorem
for the displayed family does not depend on that extension.

The quotient mechanism is more general than this witness, but positivity is
not automatic.  For a nilpotent Krein-self-adjoint charge $Q$, the finite
Kugo--Ojima theorem identifies $\operatorname{range}Q$ as the radical of the
form restricted to $\ker Q$ and gives a nondegenerate induced form on
$\ker Q/\operatorname{range}Q$ (\Mtag).  To make that quotient positive one
must additionally prove nonnegativity on $\ker Q$ and that every zero-norm
representative lies in $\operatorname{range}Q$.  The project contains an exact
three-dimensional positive witness satisfying both conditions and a matched
no-go with the same $Q$ but one ambient sign reversed: the quotient remains
nondegenerate while its surviving class has norm $-1$ (\Mtag).  Hence
Krein-self-adjointness is stable algebraic infrastructure, not a hidden
positivity assumption; the inertia and constraint hypotheses are load-bearing.

There is therefore no canonical notation $(\ker Q/\operatorname{range}Q)^+$.
The hypothesis $Q^\#=Q$ makes the Krein form descend to the quotient, but an
indefinite quotient has no preferred positive part.  A physical Hilbert space
requires additional decoder data: a chosen nonzero $D$-invariant positive
subspace of the induced-form cohomology.  Maximality is a natural further
condition; existence of a nonzero invariant choice is model-dependent, and
uniqueness is not expected.

\subsection{A positive Hodge decoder}

The quotient picture admits a sharper finite formulation.  Let $Q^*$ denote
the ordinary Hilbert adjoint of the constraint differential and define the
constraint Hodge Laplacian
\begin{equation}\label{eq:constraintHodge}
  \Delta_Q=Q^*Q+QQ^*.
\end{equation}
The generic finite Hilbert--Hodge theorem is now kernel-checked: when $Q^2=0$,
every closed vector has a unique representative in
$\ker Q\cap\ker Q^*$ modulo $\operatorname{range}Q$; a decoder commuting with
$Q$ descends to cohomology, and commuting also with $Q^*$ preserves harmonic
representatives (\texttt{GenericFiniteHodge}, \Mtag).  This theorem uses the
auxiliary positive Hilbert adjoint.  It must not be confused with the spectral
mass operator $D^\#D$ or with selection of a positive Krein subspace.

The corresponding Krein-adjoint statement is false.  On $\C^2$, take
\[
 J=\begin{pmatrix}0&1\\1&0\end{pmatrix},\qquad
 Q=\begin{pmatrix}0&1\\0&0\end{pmatrix}.
\]
Then $Q^2=0$, $Q^\#=Q$, and $Q^\#Q+QQ^\#=0$, so every vector is
Krein-harmonic; nevertheless $\ker Q=\operatorname{range}Q$, and the
cohomology is zero.  An explicit harmonic vector is not even closed
(\texttt{KreinHodgeNoGo.krein\_hodge\_no\_go}, \Mtag).

\begin{proposition}[Nonvacuous positive Hodge decoder, \Mtag]
For the explicit three-dimensional Kugo--Ojima witness, every $Q$-closed vector
is cohomologous to a $\Delta_Q$-harmonic vector.  The class represented by
$e_2$ is closed, non-exact, harmonic, and has positive Krein norm.  There is a
separate decoder $D_\mu$, commuting with $Q$ and Krein-self-adjoint, for which
\begin{equation}\label{eq:positiveHodgeMass}
  D_\mu^\#D_\mu e_2=\mu^2e_2.
\end{equation}
At $\mu=2$ this gives the explicit positive spectral cost $4$.  Reversing only
the surviving Krein sign leaves $Q$, its cohomology, and $\Delta_Q$ unchanged,
but gives $\langle e_2,e_2\rangle_J=-1$.
\end{proposition}

Thus ``particle as positive eigen-code'' now has a finite witness with all four
ingredients visible: quotient, Hilbert-harmonic representative, selected
positivity, and a nonzero spectral cost.  What remains open is derivation and
classification of nonzero $D$-invariant maximal positive sectors, derivation of
the decoder rather than insertion of $D_\mu$, comparison of inequivalent
positive choices, and stability under refinement to a physical continuum
theory.

There is also a first controlled form of decoder non-rigidity.  If $Q^2=0$,
$[D,Q]=0$, and
\begin{equation}\label{eq:chainGauge}
  D'=D+QR+RQ,
\end{equation}
then $D'$ also commutes with $Q$ and, for every closed representative $x$,
$D'x-Dx=Q(Rx)$ is exact (\Mtag).  More generally, the cross-carrier equation
$D'U-UD=Q'R+RQ$ intertwines the induced cohomology actions whenever
$UQ=Q'U$.  In the positive witness, $D'_\mu=D_\mu+2Q$ is a genuinely distinct
prephysical matrix but acts identically on the surviving class $e_2$.  This
supports a moduli interpretation of some presentation freedom; it does not show
that every alternative four-channel split is chain-homotopic.

\section{Finite quantum histories}

\subsection{The exact checkerboard sum}

Consider a $1+1$ dimensional null walk with right/left steps.  If $c(h)$ is the
number of direction changes of a history $h$, define its Gaussian-rational
amplitude by
\begin{equation}\label{eq:historyamp}
  A(h)=(\ii\epsilon m)^{c(h)}.
\end{equation}
The use of $\Q(\ii)$ makes the finite theorem exact; no floating-point or
analytic approximation enters.

\begin{theorem}[Exact finite checkerboard kernel, \Mtag]\label{thm:checker}
For a right-incoming, right-terminating path with $p$ right steps and $q>0$
left steps,
\begin{equation}\label{eq:rrkernel}
 K_{RR}(p,q)=\sum_{r=1}^{p+q}
 \binom{p}{r}\binom{q-1}{r-1}(\ii\epsilon m)^{2r}.
\end{equation}
For a right-incoming, left-terminating path,
\begin{equation}\label{eq:rlkernel}
 K_{RL}(p,q)=\sum_{r=0}^{p}
 \binom{p}{r}\binom{q-1}{r}(\ii\epsilon m)^{2r+1}.
\end{equation}
These closed forms equal the literal finite sums over histories.  The same
kernel satisfies the one-step recursion
\begin{equation}\label{eq:diracrec}
 K_{n+1}(x,d)=K_n(x-d,d)+K_n(x-d,\bar d)(\ii\epsilon m).
\end{equation}
At $m=0$, every history with a corner has zero amplitude; only straight null
histories remain.
\end{theorem}

The explicit history $[R,L,R]$ has two corners, amplitude $-1$ at
$(\epsilon,m)=(1,1)$, and amplitude $-4$ at $(1,2)$ (\Mtag).  This rules out a
vacuous path-sum statement in which the mass parameter never affects a
history.

Theorem \ref{thm:checker} is the first dynamical bridge from the null-edge
ontology to phase interference.  It does not prove a continuum propagator,
uniform convergence, or a $3+1$ dimensional checkerboard.  The existing
continuum module proves the finite symbol identities
\begin{equation}\label{eq:walksymbol}
 U_a(k)=e^{-\ii ka\sigma_z}e^{-\ii ma\sigma_x},
 \qquad \tr U_a=2\cos(ka)\cos(ma),
 \qquad \det U_a=1.
\end{equation}
Thus the exact lattice dispersion relation is
\begin{equation}\label{eq:latticedispersion}
 \cos(\omega a)=\cos(ka)\cos(ma)
 \qquad(\Mtag).
\end{equation}
At $k=0$, the two $\sigma_x$ eigenvectors have exact eigenvalues
\begin{equation}\label{eq:zeromomentumgap}
 U_a(0)x_+=e^{-\ii ma}x_+,
 \qquad
 U_a(0)x_-=e^{+\ii ma}x_-
 \qquad(\Mtag).
\end{equation}
Consequently $m$ is exactly the positive zero-momentum quasienergy on the
principal branch, while the two eigenphases are separated by $2|m|a$.  The
branch qualification is load-bearing: lattice quasienergy is periodic modulo
$2\pi/a$, so \eqref{eq:latticedispersion} does not define a globally unique
real $\omega$.

The infinitesimal generator agrees with the continuum Dirac symbol:
\begin{equation}\label{eq:generator}
 \left.\frac{\dd}{\dd\epsilon}U(k\epsilon,m\epsilon)\right|_{\epsilon=0}
 =-\ii(k\sigma_z+m\sigma_x).
\end{equation}
This leading-order statement now has an explicit finite error bound.  For a
$2\times2$ complex matrix let
$\lVert M\rVert_{\max}=\max_{i,j}|M_{ij}|$, and define
\begin{equation}\label{eq:remainderconstant}
 C(k,m)=2k^2+2m^2+|k|m^2+k^2|m|+|k||m|.
\end{equation}
Then, for $|\epsilon|\leq1$,
\begin{equation}\label{eq:remainder}
 \left\lVert U(k\epsilon,m\epsilon)
 -\bigl(1-\ii\epsilon(k\sigma_z+m\sigma_x)\bigr)\right\rVert_{\max}
 \leq C(k,m)\epsilon^2
\end{equation}
(\Mtag).  The proof is entrywise and global on the displayed interval.  It also
gives the exact massless control $U(k,0)=U_{\rm shift}(k)$ with zero
off-diagonal turn amplitude.  At $(k,m)=(3,4)$ the generator square is
$25\,I_2$ and the algebraic cross coefficient is nonzero.  Equation
\eqref{eq:remainder} is quantitative fixed-momentum, one-step control; it is
not a Trotter limit, a spacetime-walk convergence theorem, or a continuum
propagator theorem.

\subsection{A $3+1$ tetrahedral kinematic lift}

There is no single-particle obstruction to a $3+1$ checkerboard-like path
sum.  Explicit Weyl and Dirac quantum walks already provide such constructions
\cite{FosterJacobson,DarianoWeyl,NzonganiTetra}.  The first local rung needed
here can be stated over the rationals.  Let $n_0,\ldots,n_3$ be the four
vertices of a regular tetrahedron in a three-dimensional Euclidean subspace,
normalized by
\begin{equation}\label{eq:tetgram}
 n_i\mathbin{\cdot}n_j=
 \begin{cases}1,&i=j,\\-1/3,&i\ne j.\end{cases}
\end{equation}
For counts $c_i\in\N$, define the endpoint of a null history by
\begin{equation}\label{eq:tetendpoint}
 X(c)=\left(\sum_i c_i,\ \sum_i c_i n_i\right).
\end{equation}
Each primitive step $(1,n_i)$ is null, and exact expansion of
\eqref{eq:tetgram} gives
\begin{equation}\label{eq:tetmass}
 X(c)^2=\frac{8}{3}\sum_{i<j}c_i c_j\geq0
 \qquad(\Mtag).
\end{equation}
Thus every one-direction history is null, whereas the explicit history using
directions $0$ and $1$ once each has $X^2=8/3>0$.  This is a genuine $3+1$
mass-from-direction-disagreement identity.  It records endpoint kinematics;
unitarity and continuum convergence are not yet included.

The order-sensitive spin layer is now exact as well.  Let
\begin{equation}\label{eq:tetprojector}
 P_i=\frac12\bigl(I_2+r\,w_i\mathbin{\cdot}\sigma\bigr),
 \qquad 3r^2=1.
\end{equation}
The four $P_i$ are Hermitian idempotents of trace one and resolve
$\sum_iP_i=2I_2$.  Distinct directions obey the bend law
\begin{equation}\label{eq:tetbend}
 P_iP_jP_i=\frac13P_i\qquad(i\ne j),
\end{equation}
while the first oriented triple has
\begin{equation}\label{eq:tetphase}
 \tr(P_0P_1P_2)=\frac{\ii r}{3},
 \qquad
 \tr(P_2P_1P_0)=-\frac{\ii r}{3}
 \qquad(\Mtag).
\end{equation}
Chronological concatenation composes the ordered projector products in operator
order.  Hence the spin amplitude remembers bends and orientation; it is not a
function of endpoint counts alone.  This closes the first ordered-projector
rung, but no sum over all $3+1$ histories or exactly unitary massive update is
claimed.

The same tetrahedron exposes a useful normalization obstruction.  On its
three-dimensional span,
\begin{equation}\label{eq:tettightframe}
 \frac14\sum_{i=0}^3 n_i(n_i\mathbin{\cdot}v)=\frac13v
 \qquad(\Mtag).
\end{equation}
Consequently the simplest equally weighted tetrahedral projector recursion
with microscopic step scale $s$ has first-moment continuum speed $s/3$;
matching continuum $c=1$ requires $s=3$.  One therefore cannot simultaneously
demand microscopic step speed one, this normalized four-projector recursion,
and emergent speed one.  This is a no-go for that normalization, not for
$3+1$ path sums.  Foster--Jacobson make the corresponding three-times-light
step explicit, while other unitary walks use different lattice/update
structures \cite{FosterJacobson,DarianoWeyl,NzonganiTetra}.

The $D_4$ root lattice is useful here as a symmetry envelope, but not as a
drop-in null lattice.  In the integer-rescaled coordinates the four
tetrahedral steps are Hadamard vectors whose span $L_H$ lies in $D_4$ with
index eight.  This supplies an eight-sheeted Brillouin-torus cover and a
natural coset/triality bookkeeping problem.  However, the exact causal
classification is
\begin{equation}\label{eq:d4causal}
 q(h_A)=0,\qquad
 q\left(\frac{h_A+h_B}{2}\right)=-\frac23,\qquad
 q\left(\frac{h_A-h_B}{2}\right)=\frac23
 \quad(A\ne B)
 \qquad(\Mtag).
\end{equation}
Thus the extra $D_4$ root-neighbor steps include timelike and spacelike
primitives.  A separate graph theorem proves that motion using only the
$L_H$ shifts has connected components exactly the $L_H$ cosets.  The explicit
parity map
\begin{equation}\label{eq:d4quotient}
 \varphi(c_0,c_1,c_2,c_3)
  =(c_1\bmod2,\ c_3\bmod2,\ c_0+c_2\bmod2)
\end{equation}
is surjective and has kernel exactly the Hadamard image.  Therefore
\begin{equation}\label{eq:d4eight}
 D_4/L_H\simeq(\Z/2)^3,
 \qquad |D_4/L_H|=8
 \qquad(\Mtag),
\end{equation}
and the null-shift graph has exactly eight disconnected copies.  This closes
the earlier conditional asterisk.  Equation \eqref{eq:d4causal} remains the
essential warning: the full $D_4$ root-neighbor graph is not null-only, so
$D_4$ is a symmetry envelope and bookkeeping lattice rather than a license to
add every root as a microscopic propagation step.

\subsection{Null front speed across lower-order channels}

The phrase ``every particle and channel moves at $c$'' is too strong: an
interaction channel has no velocity observable, different Dirac velocity
components do not commute, and scalar and vector fields do not share that
operator.  The representation-independent statement concerns characteristics.
For a nonempty $d$-component multiplet with common wave principal symbol
\begin{equation}\label{eq:principalsymbol}
 q(\omega,k)I_d,\qquad q(\omega,k)=\omega^2-|k|^2,
\end{equation}
one has
\begin{equation}\label{eq:characteristicdet}
 \det\bigl(q(\omega,k)I_d\bigr)=q(\omega,k)^d,
 \qquad \det=0\ \Longleftrightarrow\ \omega^2=|k|^2
 \qquad(\Mtag).
\end{equation}
If aperture, closure, turn, and internal soldering mixing enter through a
finite zero-order matrix $B$, then principal scaling gives
\begin{equation}\label{eq:lowerscaling}
 \epsilon^2L(\omega/\epsilon,k/\epsilon)
 =q(\omega,k)I_d+\epsilon^2B
 \longrightarrow q(\omega,k)I_d
 \qquad(\Mtag).
\end{equation}
The first-order analogue is
$\epsilon D(p/\epsilon)=P(p)+\epsilon B$.  Hence lower-order mixing does not
change the null wavefront cone.  The hypothesis is load-bearing: an explicit
principal-order perturbation changes the determinant, and a soldering field
that changes the principal metric changes the cone rather than preserving it.

Mass affects drift without changing the front cone.  On
$E^2=|k|^2+m^2$, the group speed satisfies
\begin{equation}\label{eq:groupspeed}
 |v_g|^2=\frac{|k|^2}{E^2}=1-\frac{m^2}{E^2}<1
 \quad(m\ne0),
\end{equation}
with the exact $3$--$4$--$5$ witness $|v_g|^2=9/25$ (\Mtag).  Equations
\eqref{eq:characteristicdet}--\eqref{eq:groupspeed} therefore support the
precise claim: fields whose equations share the null principal symbol have
front velocity $c$, while massive excitations have subluminal group drift and
lower-order internal channels do not alter the front velocity.  Applying this
to every Standard Model field remains conditional on a field-by-field
principal-symbol audit.

\subsection{Action and generator boundary}

The checkerboard carrier now supplies a state-level four-channel action without
four independently chosen numbers.  For a real two-component state $\psi$, set
\begin{equation}\label{eq:channelaction}
 (S_A,S_C,S_T,S_E)=
 (\psi^TQ_A\psi,\psi^TQ_C\psi,
  4\psi^TQ_T\psi,4\psi^TE_\#\psi).
\end{equation}
The exact carrier-square identity implies
\begin{equation}\label{eq:actionbudget}
 S_A+S_C+S_T+S_E
 =\psi^T(4D^\#D)\psi
\end{equation}
for every $\psi$ (\Mtag).  At the on-shell symbol
$(E,k,m)=(5,3,4)$ and $\psi=(1,0)$, the tuple is
$(64,0,64,0)$ and the total is $128$ (\Mtag).  Thus the state-level action
budget is carrier-derived and nonvacuous.

A complementary theorem makes the phase bookkeeping local on histories.  If
each event $a$ has channel label $\chi(a)\in\{A,C,T,E\}$ and the four real
weights are $w_X$, define
\begin{equation}\label{eq:historyaction}
  S_w(h)=\sum_{a\in h}w_{\chi(a)},
  \qquad \Phi_w(h)=\exp(\ii S_w(h)).
\end{equation}
Then
\begin{equation}\label{eq:historycompose}
  S_w(h_1h_2)=S_w(h_1)+S_w(h_2),\qquad
  \Phi_w(h_1h_2)=\Phi_w(h_1)\Phi_w(h_2)
\end{equation}
(\Mtag), and the action decomposes exactly by the four event counts.  The
recursive flat-checkerboard encoding has one aperture event per null step,
one turn event per direction reversal, and no closure or soldering events.
With $(w_A,w_C,w_T,w_E)=(0,0,\pi/2,0)$,
\begin{equation}\label{eq:localcornerphase}
  \Phi_w(h)=\ii^{c(h)},\qquad
  (\epsilon m)^{c(h)}\Phi_w(h)=(\ii\epsilon m)^{c(h)}
\end{equation}
(\Mtag), including an explicit nonzero one-turn witness.  Thus the
checkerboard weight is exactly a local action phase, not merely a global
corner-count convention.

Closure now has a separate history-local phase rather than only an event
counter.  For complex edge transport $U_{xy}$, let
\begin{equation}\label{eq:closureholonomy}
  {\cal H}_U(x_0,\ldots,x_n)=
    U_{x_0x_1}U_{x_1x_2}\cdots U_{x_{n-1}x_n}.
\end{equation}
Under a unit-modulus vertex gauge $g$,
$U_{xy}'=g_yU_{xy}\overline{g_x}$, exact telescoping gives
\begin{equation}\label{eq:closuregauge}
  {\cal H}_{U'}=g_{x_n}{\cal H}_U\overline{g_{x_0}},
  \qquad x_n=x_0\Longrightarrow {\cal H}_{U'}={\cal H}_U
  \qquad(\Mtag).
\end{equation}
Holonomies multiply under history concatenation.  An exact four-vertex loop
has ${\cal H}_U=\ii\ne1$ and remains $\ii$ under a nonidentity gauge
transformation (\Mtag), so the closure phase is nonvacuous.

The four general action weights remain inputs.  The turn phase is physically
wired to the flat checkerboard and the closure phase is now wired to finite
$U(1)$ transport, but their common derivation from the carrier action is not
yet local on histories.  Nonabelian closure transport and nonzero local
soldering history densities remain open.

Finite variational modules prove that stationary pairs of a multiplier action
are exactly primal and adjoint carrier solutions, and that a constrained quadratic
action gives an eigenvalue equation (\Mtag).  Separately, a chosen Hermitian
matrix generates a unitary exponential flow preserving norm and commuting
observables (\Mtag).  These facts coexist, but they do not select $D$, $D^\#D$,
or $B(\lambda,\kappa)$ as the physical Hamiltonian.  The generator choice is a
named boundary, not an implicit theorem.

There is nevertheless an exact finite generator-to-update bridge.  For a
chosen matrix $H$, the Cayley step
\begin{equation}\label{eq:cayley}
 C_{\Delta t}(H)=
 (1+\ii\Delta t H/2)^{-1}(1-\ii\Delta t H/2)
\end{equation}
solves the cleared-denominator Crank--Nicolson equation; for Hermitian $H$
under the displayed inverse identities it is unitary (\Mtag).  The exact
fixture $H=\sigma_x$, $\Delta t=2$ gives
$C_2(\sigma_x)=-\ii\sigma_x$ and $e_0\mapsto-\ii e_1$ (\Mtag).  This derives
the update from a supplied generator, not the physical generator itself.

\section{Closure, binding, and finite many-body structure}

The closure channel is signed.  On the explicit physical quotient its induced
form has balanced inertia $(2,2,0)$ (\Mtag).  Consequently it cannot be treated
as a positive energy density.  The positive Wilson/plaquette closure-defect
energy is a different quadratic object.  This distinction prevents a common
but false identification of a chromomagnetic term with total gluon energy.

The finite fermionic occupation model has a sharp free result: if the
one-particle energies have least value $\lambda-\kappa>0$, then the free
second-quantized gap above the vacuum is exactly $\lambda-\kappa$ (\Mtag).
Free two-particle energies are constituent sums and do not bind.

A separate $\Lambda^2$ interaction is derived by second-quantizing a closure
curvature.  For a sorted three-mode spectrum and $\kappa>0$, closure acting in
the excited-mode plane gives a least eigenvalue strictly below the free
two-particle threshold (\Mtag).  Closure in a plane containing the ground mode
instead leaves the lowest pair unbound whenever
\begin{equation}\label{eq:wrongplane}
 \kappa^2\leq(d_2-d_0)(d_2-d_1).
\end{equation}
Thus ``closure binds'' is not a theorem about arbitrary curvature placement.
The actual finite carrier curvature is machine-identified with the binding
plane, closing that choice for the witness.  A continuum QCD or confinement
interpretation remains \Ctag.

\section{Protected masslessness and family structure}

For a finite odd operator between positive and negative chiral sectors,
rank--nullity gives
\begin{equation}\label{eq:index}
 \dim\ker D\geq n_+-n_-.
\end{equation}
The lower bound survives odd mass perturbations (\Mtag).  A separate finite
winding model realizes index $w$, kernel dimension $w$, and refinement-stable
protected mode count (\Mtag).  These are finite index statements, not a
continuum anomaly theorem.

The project also formalizes the Kobayashi--Maskawa phase-counting spine.  Two
families admit no physical CP phase after rephasing, while three families admit
an explicit nonzero rational Jarlskog witness (\Mtag).  What does not follow is
equally important: triality, anomaly cancellation, and the finite carrier
category do not by themselves force exactly three replicated charge families
(\Mtag no-go).  More concretely, diagonal triplication retains the full family
permutation commutant, and no family label is fixed by every permutation
(\Mtag).  Three copies therefore do not become three canonically ordered
generations.  A rank-fixing and symmetry-breaking input is still required.  The finite theory
therefore reproduces the special role of three in CP phase counting without
claiming to derive the observed number of generations.

\section{Conjecture scoreboard}

The physical dictionary is useful only while its failure conditions remain
visible.  Table~\ref{tab:scoreboard} separates the proposed identification from
the finite result that currently supports it.

\begin{table}[H]
\centering
\small
\begin{tabularx}{\linewidth}{@{}>{\raggedright\arraybackslash}p{0.22\linewidth}>{\raggedright\arraybackslash}p{0.36\linewidth}X@{}}
\toprule
Conjectural identification & Kill condition & Current status\\
\midrule
Null histories recover free relativistic propagation & No controlled
many-step convergence to the Dirac propagator & Exact finite path sum,
recursion, dispersion, eigenphases, and one-step $O(\epsilon^2)$ bound; the
many-step theorem remains open\\
The four carrier blocks match continuum field sectors &
Refinement converges to different operators or requires inserted projectors &
Concrete four-block coefficient rigidity, but abstract non-uniqueness\\
Closure supplies a QCD-like binding channel & Singlet low modes disappear
under enlargement/refinement, or the physical curvature occupies the no-binding
plane & Exact finite binding and wrong-plane control; no confinement theorem\\
Soldering supplies gravitational geometry & No nondegenerate continuum coframe,
covariant action, or accepted field equation emerges & Finite nondegenerate
matrix coframe with metric, volume, covariance, refinement, and nonzero action\\
Finite family structure selects three generations & Replication symmetry leaves
family labels interchangeable, or another rank is equally admissible & Exact
three-family CP witness, together with rank-fixing and permutation no-go theorems\\
$D_4$ is the microscopic step graph & Additional roots are not null &
$D_4/L_H\simeq(\Z/2)^3$ is now exact, but causal classification restricts
primitive propagation to the tetrahedral Hadamard sublattice\\
Particles are positive Hodge eigen-codes & The physical quotient is empty or
indefinite, or the spectral decoder fails to descend under refinement & Exact
generic Hilbert--Hodge representatives; exact positive class with spectral
cost $4$; matched sign no-go; and a $2\times2$ Krein-Hodge collapse.  Nonzero
invariant positive-sector selection and continuum reconstruction remain open\\
\bottomrule
\end{tabularx}
\caption{Current claim/kill/status scoreboard.  Every physical identification
in the first column remains \Ctag even when the supporting finite statement is
\Mtag.}
\label{tab:scoreboard}
\end{table}

\section{Falsifiers and claim boundaries}
\label{sec:falsifiers}

The model is intended to become easier to reject as it grows.  Its current
boundaries can be stated as concrete tests.

\begin{enumerate}[leftmargin=1.7em]
  \item \textbf{Continuum dictionary.}  If refined aperture, closure, turn, and
  soldering blocks do not converge to kinetic, gauge-curvature, scalar, and
  geometric operators, the physical names must be retired in favor of neutral
  block labels.
  \item \textbf{Carrier rigidity.}  The explicit witness is coefficient-rigid,
  but the category is not.  A proposed extra principle must remove known
  inequivalent complements without simply encoding the desired carrier.
  \item \textbf{Path dynamics.}  The finite checkerboard theorem must extend to
  a controlled continuum limit.  Failure to reproduce the free Dirac
  propagator would kill the path interpretation.
  \item \textbf{External dispersion gate.}  The exact principal-branch relation
  \eqref{eq:latticedispersion} has the small-$a$ massive expansion
  \[
    \omega^2=k^2+m^2-\frac{a^2}{3}k^2m^2+O(a^4).
  \]
  Precision bounds on species-dependent massive-particle dispersion therefore
  bound any literal microscopic spacing $a$ (\Ttag/\Ctag).  In the minimal
  massless sector, however, $m=0$ gives $\omega=|k|$ exactly on the principal
  branch.  Gamma-ray photon time-of-flight is consequently not the leading
  constraint on this particular walk; importing such a bound without a photon
  correction would be spurious.
  \item \textbf{Resource interpretation.}  Determinant/linear entropy must be
  monotone only under a precisely stated allowed-operation class.  A class that
  includes the known signed-closure counterexample is false.
  \item \textbf{Structured closure.}  Robust singlet low modes must persist
  beyond the current small matrices and under refinement.  Random disorder is
  already known not to supply the desired mechanism.
  \item \textbf{Gravity.}  The soldering-gradient term must obey a legitimate
  local-frame transformation law and approach a known gravitational action.
  The rational matrix-coframe defect now passes finite nondegeneracy, induced
  metric, volume, Lorentz-covariance, invariant-action, and refinement tests,
  but a nondegenerate continuum tetrad field and a gravitational limit remain
  unproved.  The present finite
  torsion/nonmetricity split is not continuum general relativity.
\end{enumerate}

No result in this paper computes an electron, quark, proton, or cosmological
mass.  No theorem establishes a continuum Standard Model, a physical path
measure, scattering amplitudes, or quantum gravity.  The strongest result is a
coherent finite chain: canonical null-edge mass kinematics, a positive carrier
witness, a rigid concrete block presentation, an exact finite history sum, and
an exact quasienergy spectrum tying the corner parameter to the zero-momentum
walk gap.

\section{Machine verification and reproducibility}

The Lean kernel is the proof checker.  The project uses Lean 4.28.0 and a pinned
Mathlib snapshot.  Headline draft results include build-enforced
\texttt{\#print axioms} message guards; the expected footprint is limited to
standard quotient/extensionality/classical-choice principles.  Numerical or
computer-algebra scripts are treated as external oracles and are not used as
proofs.

For example, the checked guard output for the new dispersion and concrete
$D_4$ verdicts is, verbatim apart from line wrapping,
\begin{verbatim}
depends on axioms: [propext, Classical.choice, Quot.sound]
\end{verbatim}
and the build fails if that footprint changes.

\clearpage
\begin{landscape}
\begin{table}[!ht]
\centering
\scriptsize
\begin{tabularx}{\linewidth}{@{}>{\raggedright\arraybackslash}p{0.34\linewidth}>{\raggedright\arraybackslash}p{0.35\linewidth}X@{}}
\toprule
Lean declaration & Source module & Mathematical payload\\
\midrule
\path{i1_5_cauchy_binet_mass_identity} &
\path{NullEdge/GateI1/Core.lean} & Equation \eqref{eq:cb}\\
\path{spinorWedge_sl2_invariant} &
\path{PhysicsSM/Spinor/PluckerMassCovariance.lean} & Equation \eqref{eq:termwiseSL2}\\
\path{a2_sqrt_minkowskiSq_add_ge_of_futureCone} &
\path{NullEdge/GateI1/Core.lean} & Free mass superadditivity \eqref{eq:superadditive}\\
\path{four_mul_det_gram_eq_concurrence_sq} &
\path{NullEdge/TwoEdgeMassConcurrence.lean} & Concurrence identity \eqref{eq:concurrence}\\
\path{compositeMassSq_eq_zero_iff_collinear} &
\path{GateI1/CompositeApertureMass.lean} & Massless iff common null direction\\
\path{detP_unique} &
\path{NullEdge/DetPUniqueness.lean} & Theorem \ref{thm:detunique}\\
\path{free_mass_operator_eq_plucker} &
\path{NullEdge/Carrier/FreeMassBridge.lean} & Free spectral/Gram bridge\\
\path{path_mass_visibility_duality} &
\path{NullEdge/MassCoherencePathEquivalence.lean} & Equation \eqref{eq:duality}\\
\path{mass_monotone_under_pinch} &
\path{NullEdge/EntropyMonotoneReal.lean} & Equation \eqref{eq:pinch}\\
\path{carrier_krein_square} &
\path{Carrier/CarrierKreinSquare.lean} & Equation \eqref{eq:master}\\
\path{four_channels_linearIndependent} &
\path{NullEdge/FourChannelRigidityCapstone.lean} & Concrete coefficient rigidity\\
\path{T2_positive_mass} &
\path{NullEdge/Carrier/SectorGroundMassWitness.lean} & Explicit positive physical sector\\
\path{nonvacuous_positive_sector}, \path{nondegenerate_but_indefinite_no_go} &
\path{NullEdge/Carrier/KreinPositiveSectorWitness.lean} & Positive quotient mechanism and matched sign no-go\\
\path{closed_cohomologous_to_harmonic}, \path{positive_hodge_mass_witness} &
\path{NullEdge/Carrier/PositiveHodgeDecoder.lean} & Equations \eqref{eq:constraintHodge}--\eqref{eq:positiveHodgeMass}\\
\path{chainHomotopy_physical_action}, \path{decoder_chain_intertwining_packet} &
\path{NullEdge/Carrier/DecoderChainHomotopy.lean} & Channel-gauge equivalence \eqref{eq:chainGauge}\\
\path{B_posDef_iff}, \path{B_least_eigenvalue} &
\path{Carrier/MassGapWitness.lean} & Exact phase diagram \eqref{eq:B}\\
\path{exact_path_sum_eq_closed_kernel} &
\path{NullEdge/ExactCheckerboardPathSum.lean} & Equation \eqref{eq:rrkernel}\\
\path{exact_path_sum_dirac_recursion} &
\path{NullEdge/ExactCheckerboardPathSum.lean} & Equation \eqref{eq:diracrec}\\
\path{Ustep_hasDerivAt_generator} &
\path{Carrier/ContinuumLimit.lean} & Equation \eqref{eq:generator}\\
\path{Ustep_taylor_remainder_bound} &
\path{NullEdge/QuantitativeDiracWalkContinuum.lean} & Equation \eqref{eq:remainder}\\
\path{exact_quantum_walk_dispersion_verdict} &
\path{NullEdge/ExactQuantumWalkDispersion.lean} & Equations \eqref{eq:latticedispersion}--\eqref{eq:zeromomentumgap}\\
\path{endpointMassSq_eq_pairDisagreement} &
\path{NullEdge/TetrahedralNullHistory.lean} & Equation \eqref{eq:tetmass}\\
\path{frameOp_eq_third} &
\path{NullEdge/TetrahedralNullHistory.lean} & Equation \eqref{eq:tettightframe}\\
\path{tetrahedral_spin_projector_path_verdict} &
\path{NullEdge/TetrahedralSpinProjectorPath.lean} & Equations \eqref{eq:tetprojector}--\eqref{eq:tetphase}\\
\path{d4_added_step_causal_verdict} &
\path{NullEdge/GateC1/D4AddedStepCausalClassification.lean} & Equation \eqref{eq:d4causal}\\
\path{componentEquiv}, \path{eightCopies} &
\path{NullEdge/GateC1/D4DisconnectedCopy.lean} & Coset-component theorem\\
\path{d4_concrete_quotient_verdict} &
\path{NullEdge/GateC1/D4ConcreteQuotient.lean} & Equations \eqref{eq:d4quotient}--\eqref{eq:d4eight}\\
\path{principal_characteristic_iff_null} &
\path{NullEdge/LowerOrderChannelCausality.lean} & Equations \eqref{eq:principalsymbol}--\eqref{eq:characteristicdet}\\
\path{four_lowerOrder_channels_preserve_principal_symbol} &
\path{NullEdge/LowerOrderChannelCausality.lean} & Equation \eqref{eq:lowerscaling}\\
\path{massive_groupSpeedSq_lt_one} &
\path{NullEdge/LowerOrderChannelCausality.lean} & Equation \eqref{eq:groupspeed}\\
\path{checkerboard_channel_action_total} &
\path{NullEdge/FourChannelPathActionCapstone.lean} & Equation \eqref{eq:actionbudget}\\
\path{checkerboardAmplitude_eq_corner_power} &
\path{NullEdge/HistoryLocalFourChannelAction.lean} & Equations \eqref{eq:historyaction}--\eqref{eq:localcornerphase}\\
\path{u1_history_closure_holonomy_verdict} &
\path{NullEdge/U1HistoryClosureHolonomy.lean} & Equations \eqref{eq:closureholonomy}--\eqref{eq:closuregauge}\\
\path{cayley_hamiltonian_generator_verdict} &
\path{NullEdge/CayleyHamiltonianGenerator.lean} & Equation \eqref{eq:cayley}\\
\path{soldering_local_frame_covariance_verdict} &
\path{NullEdge/SolderingLocalFrameCovariance.lean} & Equations \eqref{eq:soldercov}--\eqref{eq:solderrefine}\\
\path{nondegenerate_soldering_geometry_verdict} &
\path{NullEdge/NondegenerateSolderingGeometry.lean} & Equations \eqref{eq:inducedmetric}--\eqref{eq:rationalboost}\\
\path{secondQuantized_massGap} &
\path{Carrier/FockMassGap.lean} & Free finite Fock gap\\
\path{derived_boundState_below_threshold} &
\path{Carrier/DerivedInteraction.lean} & Closure-derived binding\\
\path{derived_wrongPlane_no_binding} &
\path{Carrier/DerivedInteraction.lean} & Obstruction \eqref{eq:wrongplane}\\
\path{winding_protects_low_modes} &
\path{NullEdge/WindingLowModes.lean} & Finite index protection\\
\path{jarlskog_Vwitness_ne_zero} &
\path{NullEdge/FiniteKMCP.lean} & Nonzero finite CP witness\\
\path{three_generations_not_forced} &
\path{NullEdge/FamilyRankNoGo.lean} & Generation-count no-go\\
\path{generation_permutation_nogo} &
\path{NullEdge/GenerationPermutationNoGo.lean} & Diagonal-triplication symmetry no-go\\
\path{one_pair_joint_chain} &
\path{NullEdge/Carrier/PluckerJointTheoryWitness.lean} & One-pair mass/action/flow/ensemble composition\\
\path{one_pair_drives_dirac_symbol} &
\path{NullEdge/Carrier/PluckerDiracCarrierBridge.lean} & Same pair mass in Gibbs and internal $3+1$ Dirac symbol\\
\path{axisBasis_conjugates_velocity} &
\path{NullEdge/CliffordDiagonalPositionBridge.lean} & Explicit component-sign-to-Clifford dictionary\\
\path{pathAction_eq_zero_iff_evolution} &
\path{NullEdge/FiniteUnitaryPathAction.lean} & Zero residual action iff selected unitary EOM\\
\path{outcome_prob_sum_one} &
\path{NullEdge/FiniteInstrumentAPI.lean} & Kraus-complete outcome normalization with imported trace rule\\
\bottomrule
\end{tabularx}
\caption{Selected machine-verified anchors.  Paths are relative to
\texttt{PhysicsSM/Draft} unless the full \texttt{PhysicsSM/} prefix is shown.}
\label{tab:anchors}
\end{table}
\end{landscape}
\clearpage

No theorem cited in Table~\ref{tab:anchors} is accepted merely because a
numerical fixture passes.  Capstone conjunctions are integration interfaces;
the scientific citations are the payload declarations beneath them.

This working revision is based on repository commit
\path{10a9505dffb1decfd8cfadb8e37684609acc1d8d}, with additional manuscript and
Lean changes still present in the working tree.  It is therefore not yet an
archival identifier for the exact paper state.  A clean commit, source archive,
and Zenodo DOI are explicit release gates before submission; no DOI is claimed
for this draft.

The principal verification commands are
\begin{verbatim}
lake env lean PhysicsSM/Draft/NullEdge/ExactCheckerboardPathSum.lean
lake build PhysicsSM.Draft.NullEdge.QuantitativeDiracWalkContinuum
lake build PhysicsSM.Draft.NullEdge.ExactQuantumWalkDispersion
lake build PhysicsSM.Draft.NullEdge.TetrahedralNullHistory
lake build PhysicsSM.Draft.NullEdge.TetrahedralSpinProjectorPath
lake build PhysicsSM.Draft.NullEdge.LowerOrderChannelCausality
lake build PhysicsSM.Draft.NullEdge.GateC1.D4AddedStepCausalClassification
lake build PhysicsSM.Draft.NullEdge.GateC1.D4DisconnectedCopy
lake build PhysicsSM.Draft.NullEdge.GateC1.D4ConcreteQuotient
lake build PhysicsSM.Draft.NullEdge.FourChannelRigidityCapstone
lake build PhysicsSM.Draft.NullEdge.FiniteHamiltonianGeneratorCapstone
lake build PhysicsSM.Draft.NullEdge.CayleyHamiltonianGenerator
lake build PhysicsSM.Draft.NullEdge.HistoryLocalFourChannelAction
lake build PhysicsSM.Draft.NullEdge.U1HistoryClosureHolonomy
lake build PhysicsSM.Draft.NullEdge.SolderingLocalFrameCovariance
lake build PhysicsSM.Draft.NullEdge.NondegenerateSolderingGeometry
lake build PhysicsSM.Draft.NullEdge.Carrier.GenericFiniteHodge
lake build PhysicsSM.Draft.NullEdge.Carrier.KreinHodgeNoGo
lake build PhysicsSM.Draft.NullEdge.SelfConsistentDecoder
lake build PhysicsSM.Draft.NullEdge.OvernightTheoryAxiomGuard
\end{verbatim}
A full draft build additionally depends on optional Sphere-Packing-Lean modules;
native Windows verification of the null-edge targets is independent of those
optional modules.

\section{Conclusion}

The finite theorem is simple to state.  A bundle of future null spinors defines
a positive momentum Gram matrix, and its determinant is the sum of squared
pairwise wedges.  Rank detects the massless/massive dichotomy; the determinant
is the continuous spectral area of mutually disagreeing null directions and is
canonical among quadratic forms vanishing on the full null boundary.

The carrier model adds structure without changing that invariant.  Its square
separates four operator types; the explicit rational presentation is
coefficient-rigid and admits a positive sector.  Closure can lower the finite
gap and, in the correct two-body plane, bind below threshold.  Turn mass enters
an exact checkerboard history only at corners; the corner weight is now the
exponential of an additive local action, producing a genuine finite phase sum
and discrete Dirac recursion.  The same walk has exact dispersion
$\cos(\omega a)=\cos(ka)\cos(ma)$ and zero-momentum eigenphases
$e^{\mp\ii ma}$, so ``spectral'' and ``turn mass'' meet in one checked finite
object.  Closure separately carries a nontrivial
history-local $U(1)$ holonomy that is gauge invariant on closed loops.  A first
$3+1$ rational tetrahedral lift now
proves null primitive steps, a positive timelike endpoint from mixed
directions, and the exact one-third normalization obstruction.  Its Hadamard
sublattice has the concrete eight-element quotient
$D_4/L_H\simeq(\Z/2)^3$, while the causal audit excludes the remaining roots as
primitive null steps.  Independently,
a principal-symbol theorem identifies the common null front cone preserved by
lower-order channel mixing while retaining subluminal massive drift.  The
soldering defect now has a nondegenerate matrix-coframe realization with an
induced Lorentz metric, determinant volume, invariant quadratic action, exact
refinement, and a nonzero rational boost witness.  A soldering field that
modifies the principal metric must nevertheless be audited as cone-defining
rather than lower order.

These statements are enough to define a serious finite research program, but
not enough to identify the model with nature.  The decisive next results are
a compact-box $3+1$ product-limit bound for the landed successive-axis walk, a
field-valued action whose reductions yield both the Pluecker recurrence and
the Dirac update, a physical position/PDE continuum lift, a derived choice of
time generator and spatial frame, and a field-by-field principal-symbol audit
of the four channels.  The simultaneous six-channel scalar-square route has
been killed rather than silently reinterpreted; Route B is a distinct
construction.  The higher-dimensional many-particle QCA obstruction of
Mlodinow--Brun is a separate locality problem, not a no-go for the landed
single-particle kinematics \cite{MlodinowBrunNoGo}.  The failure conditions are
already visible.

\begin{thebibliography}{99}

\bibitem{ElvangHuang}
H.~Elvang and Y.-t.~Huang,
\emph{Scattering Amplitudes}, arXiv:1308.1697.

\bibitem{ArkaniHamedMassive}
N.~Arkani-Hamed, T.-C.~Huang, and Y.-t.~Huang,
\emph{Scattering Amplitudes For All Masses and Spins},
JHEP 11 (2021) 070, arXiv:1709.04891.

\bibitem{ChinLee}
S.~Chin and S.~Lee,
\emph{Momentum bispinor, two-qubit entanglement and twistor space},
arXiv:1407.2492.

\bibitem{Wootters}
W.~K.~Wootters,
\emph{Entanglement of formation of an arbitrary state of two qubits},
Phys. Rev. Lett. 80 (1998) 2245--2248, arXiv:quant-ph/9709029.

\bibitem{DarianoPerinottiPrinciples}
G.~M.~D'Ariano and P.~Perinotti,
\emph{Derivation of the Dirac equation from principles of information
processing}, Phys. Rev. A 90 (2014) 062106, arXiv:1306.1934.

\bibitem{BisioDarianoTosini}
A.~Bisio, G.~M.~D'Ariano, and A.~Tosini,
\emph{Quantum field as a quantum cellular automaton: the Dirac free evolution
in one dimension}, Ann. Phys. 354 (2015) 244, arXiv:1212.2839.

\bibitem{DarianoPathIntegral}
G.~M.~D'Ariano, N.~Mosco, P.~Perinotti, and A.~Tosini,
\emph{Path-integral solution of the one-dimensional Dirac quantum cellular
automaton}, Phys. Lett. A 378 (2014) 3165--3168, arXiv:1406.1021.

\bibitem{FeynmanHibbs}
R.~P.~Feynman and A.~R.~Hibbs,
\emph{Quantum Mechanics and Path Integrals}, McGraw--Hill, 1965.

\bibitem{JacobsonSchulman}
T.~Jacobson and L.~S.~Schulman,
\emph{Quantum stochastics: the passage from a relativistic to a
non-relativistic path integral}, J. Phys. A 17 (1984) 375.

\bibitem{FosterJacobson}
B.~Z.~Foster and T.~Jacobson,
\emph{Spin on a 4D Feynman Checkerboard}, arXiv:1610.01142.

\bibitem{DarianoWeyl}
G.~M.~D'Ariano, N.~Mosco, P.~Perinotti, and A.~Tosini,
\emph{Path-sum solution of the Weyl Quantum Walk in 3+1 dimensions},
Phil. Trans. R. Soc. A 375 (2017) 20160394, arXiv:1705.08552.

\bibitem{NzonganiTetra}
U.~Nzongani, N.~Eon, I.~Marquez-Martin, A.~Perez, G.~Di Molfetta, and
P.~Arrighi,
\emph{Dirac quantum walk on tetrahedra}, arXiv:2404.09840.

\bibitem{MlodinowBrun}
L.~Mlodinow and T.~A.~Brun,
\emph{Discrete spacetime, quantum walks, and relativistic wave equations},
Phys. Rev. A 97 (2018) 042131, arXiv:1802.03910.

\bibitem{MlodinowBrunNoGo}
L.~Mlodinow and T.~A.~Brun,
\emph{Quantum field theory from a quantum cellular automaton in one spatial
dimension and a no-go theorem in higher dimensions},
Phys. Rev. A 102 (2020) 042211, arXiv:2006.08927.

\bibitem{Barrett}
J.~W.~Barrett,
\emph{A Lorentzian version of the non-commutative geometry of the Standard
Model}, J. Math. Phys. 48 (2007) 012303, arXiv:hep-th/0608221.

\bibitem{vanDenDungen}
K.~van den Dungen,
\emph{Krein spectral triples and the fermionic action},
Math. Phys. Anal. Geom. 19 (2016) 4, arXiv:1505.01939.

\bibitem{Wilson}
K.~G.~Wilson,
\emph{Confinement of quarks}, Phys. Rev. D 10 (1974) 2445.

\bibitem{BanksCasher}
T.~Banks and A.~Casher,
\emph{Chiral symmetry breaking in confining theories},
Nucl. Phys. B 169 (1980) 103.

\bibitem{GinspargWilson}
P.~H.~Ginsparg and K.~G.~Wilson,
\emph{A remnant of chiral symmetry on the lattice},
Phys. Rev. D 25 (1982) 2649.

\bibitem{NielsenNinomiya}
H.~B.~Nielsen and M.~Ninomiya,
\emph{Absence of neutrinos on a lattice}, Nucl. Phys. B 185 (1981) 20.

\bibitem{Regge}
T.~Regge,
\emph{General relativity without coordinates}, Nuovo Cim. 19 (1961) 558.

\bibitem{DittrichRyan}
B.~Dittrich and J.~P.~Ryan,
\emph{Phase space descriptions for simplicial 4d geometries},
Class. Quant. Grav. 28 (2011) 065006, arXiv:0807.2806.

\bibitem{FreidelSpeziale}
L.~Freidel and S.~Speziale,
\emph{Twisted geometries: A geometric parametrisation of $SU(2)$ phase space},
Phys. Rev. D 82 (2010) 084040, arXiv:1001.2748.

\bibitem{BaezWise}
J.~C.~Baez and D.~K.~Wise,
\emph{Teleparallel gravity as a higher gauge theory},
Commun. Math. Phys. 333 (2015) 153--186, arXiv:1204.4339.

\bibitem{ArrighiCurvedWalk}
P.~Arrighi, S.~Facchini, and M.~Forets,
\emph{Quantum walking in curved spacetime}, arXiv:1505.07023.

\bibitem{PhysLean}
J.~Tooby-Smith,
\emph{HepLean: Digitalising high energy physics}, arXiv:2405.08863.

\end{thebibliography}

\end{document}
